\documentclass[10pt,letterpaper]{article}

\usepackage[margin=0.7in]{geometry}
\usepackage[T1]{fontenc}
\usepackage{lmodern}
\usepackage{microtype}
\usepackage[hidelinks]{hyperref}
\usepackage{enumitem}
\usepackage{tabularx}
\usepackage{xcolor}
\usepackage{titlesec}
\usepackage{array}
\usepackage{fancyhdr}
\usepackage{ragged2e}
\usepackage{fontawesome5}

\setlength{\parindent}{0pt}
\setlength{\parskip}{0pt}
\setlength{\tabcolsep}{0pt}
\setlist[itemize]{leftmargin=1.1em,itemsep=1.5pt,topsep=2pt,parsep=0pt,partopsep=0pt}

% --- Palette -------------------------------------------------------------
% accent  : section labels / name
% rulegray: hairline rules
% linkblue: hyperlinks
\providecommand{\accenthex}{1A3D5C}   % deep slate blue
\definecolor{accent}{HTML}{\accenthex}
\definecolor{rulegray}{HTML}{C8CDD4}
\definecolor{mutedgray}{HTML}{55606B}
\definecolor{linkblue}{HTML}{2F5D9A}
\hypersetup{
  colorlinks=true,
  urlcolor=linkblue,
  linkcolor=linkblue,
  pdfborderstyle={/S/U/W 0}
}

% --- Section headers: small-caps label with a full-width hairline rule ----
\titleformat{\section}
  {\normalfont\scshape\small\bfseries\color{accent}}
  {}{0em}{}[\vspace{2pt}{\color{rulegray}\titlerule[0.6pt]}]
\titlespacing*{\section}{0pt}{9pt}{4pt}

% --- Experience / project entry: bold title + right-aligned date ----------
\newcommand{\entryheader}[4]{%
  \begin{tabular*}{\textwidth}{@{\extracolsep{\fill}}lr}
    \textbf{#1} & \textbf{\color{mutedgray}#2} \\
    {\color{mutedgray}\itshape #3} & {\color{mutedgray}\itshape #4} \\
  \end{tabular*}\vspace{0.18em}
}

% --- Publication entry: year column + reference -----------------------------
\newcommand{\pubentry}[2]{%
  \begin{tabularx}{\textwidth}{@{}>{\color{mutedgray}}p{0.95cm}X@{}}
    #1 & #2 \\
  \end{tabularx}\vspace{0.22em}
}

% --- Skills row -----------------------------------------------------------
\newcommand{\skillrow}[2]{\textbf{\color{accent}#1:}~~#2\par\vspace{1pt}}

% --- Highlights / metrics strip (boxed key facts) -------------------------
% Usage: \highlightstrip{item a \highlightsep item b \highlightsep ...}
\newcommand{\highlightsep}{\quad\textcolor{rulegray}{\textbar}\quad}
\newcommand{\highlightstrip}[1]{%
  \begingroup
  \setlength{\fboxsep}{7pt}%
  \noindent\colorbox{accent!6}{%
    \parbox{\dimexpr\textwidth-2\fboxsep\relax}{%
      \footnotesize\color{mutedgray}#1%
    }%
  }\par\vspace{8pt}%
  \endgroup
}

% --- Header: name, tagline, contact line ----------------------------------
% #1 = name, #2 = first contact line, #3 = second contact line
\newcommand{\makeprofileheader}[3]{%
  {\fontsize{22}{24}\selectfont\bfseries\color{accent}#1}\par
  \vspace{5pt}
  {\color{rulegray}\hrule height 1pt}
  \vspace{7pt}
  {\footnotesize\color{mutedgray}#2\par}
  \vspace{2pt}
  {\footnotesize\color{mutedgray}#3\par}
  \vspace{6pt}
}

% --- Footer ---------------------------------------------------------------
\newcommand{\setupcvfooter}[1]{%
  \pagestyle{fancy}
  \fancyhf{}
  \fancyfoot[L]{\footnotesize\color{mutedgray} Rajeev Jain \textbullet{} CV \textbullet{} Updated #1}
  \fancyfoot[C]{}
  \fancyfoot[R]{\footnotesize\color{mutedgray}\thepage}
  \renewcommand{\headrulewidth}{0pt}
  \renewcommand{\footrulewidth}{0pt}
  \setlength{\footskip}{18pt}
}

\pagestyle{empty}

\begin{document}
\sloppy

\makeprofileheader
  {Rajeev Jain}
  {\faGlobe~\href{https://rajeeja.github.io}{rajeeja.github.io} \quad \faGithub~\href{https://github.com/rajeeja}{rajeeja} \quad \faLinkedin~\href{https://linkedin.com/in/rajeeja}{rajeeja} \quad \faGraduationCap~\href{https://scholar.google.com/citations?user=bC77n9MAAAAJ}{Google Scholar} \quad \faOrcid~\href{https://orcid.org/0000-0002-1235-918X}{ORCID}}
  {\faEnvelope~\href{mailto:jain@anl.gov}{jain@anl.gov} \quad \faEnvelope~\href{mailto:rajeeja@gmail.com}{rajeeja@gmail.com} \quad \faMapMarker~Chicago, IL}

{\large\color{mutedgray} Principal Specialist, Research Software Engineering \textbar{} AI Systems and Verification \textbar{} HPC \textbar{} Scientific Computing}\par
\vspace{6pt}

\section{Profile}
Principal research software engineer, at Argonne since 2009, turning research prototypes into production systems across AI infrastructure, scientific computing, and HPC. Lead developer of UXarray, a Python library for unstructured climate-grid analysis used at DOE labs and NCAR, including an MCP server that gives AI agents typed, scoped, provenance-tracked access to Earth-system meshes on HPC clusters via Globus Compute --- the alternative to handing a team general-purpose chat and calling it a strategy. Equally focused on verifying what such systems produce: caught silent numerical drift, hidden heap allocation, and unreachable vendor throughput figures in code and benchmarks that read as correct. Ran 10,000+ cancer-AI training experiments on Summit; ported the Pangu-Weather climate emulator to Intel XPUs on Aurora. Two R\&D 100 Awards (CANDLE 2023, FLASH-X 2022).

\section{Professional Experience}
\entryheader
  {Argonne National Laboratory}
  {2009 -- Present}
  {Principal Specialist, Research Software Engineering}
  {Mathematics and Computer Science Division, Chicago, IL}
\begin{itemize}
  \item \textbf{UXarray and MCP server.} Lead developer of a Python library for unstructured climate grid analysis used at DOE labs and NCAR. Designed grid I/O for ESMF / MPAS / SCRIP / HEALPix and conservative zonal averaging. Built the UXarray MCP server: 20+ typed tools letting AI agents inspect, subset, visualize, and run analysis on Earth-system meshes locally or on HPC clusters (Argonne Improv, UCAR Derecho) via Globus Compute, with structured provenance. Presented at SciFM26.
  \item \textbf{AI code and benchmark verification.} Independent measurement applied to AI tooling and machine-assisted code. Published a reproducible harness for local LLM inference: bytes-per-token parsed from GGUF tensor tables rather than file size (which overstates 64\% on gathered-embedding models), and memory bandwidth measured at 340 GB/s against an unreachable 400 GB/s datasheet pin rate --- showing four different architectures all land within 36--41\% of bandwidth, so throughput is set by model size and runtime, not architecture. In numerical Python, traced 0.68 m of drift on a production Earth mesh to catastrophic cancellation inside a spherical cross product --- a difference of nearly equal products, not a running sum --- and proposed error-free-transformation arithmetic for the kernel (merged upstream, UXarray PR \#1513); recovered a 1.67$\times$ speedup by removing heap allocation a faithful C++-to-Python port had silently introduced.
  \item \textbf{Pangu-Weather on Aurora.} Ported the Pangu-Weather climate emulator to Aurora's Intel XPU stack (60,000+ GPUs): PMIX/PALS launcher mapping, XPU/CUDA device branching, device-aware AMP (GradScaler on CUDA, bf16 on XPU), first stable DDP baseline. Established smoke-test and queue-aware bring-up scripts separating launcher, backend, and device-placement failures. Measured one-node baseline: 12 XPU ranks under DDP, $\sim$12 s steady-state epochs after a warmup epoch 8$\times$ longer.
  \item \textbf{CANDLE / IMPROVE.} Built CANDLE/Supervisor HPO workflow (Swift/T, mlrMBO, DEAP, Hyperopt/TPE); ran 10,000+ cancer drug response experiments on Summit, Theta, and Cori. Applied noise injection and counterfactual analysis to NT3 tumor classification; identified PLOD2, RGS5, and TP53I13 as decision-boundary drivers. Led IMPROVE: \texttt{improvelib} Python package with GitHub Actions CI/CD (unit tests, Docker, Parsl CSA) and UNO dual-branch neural network for cross-study benchmarking. 15+ researchers across Argonne, LLNL, ORNL. R\&D 100 Award, 2023.
  \item \textbf{FLASH-X.} I/O and compression lead for a million-line multiphysics engine. Async HDF5 with Argobots plus SZ3 / ZFP compression reduced checkpoint overhead 40--70\% on Summit with 50\%+ storage savings. Enabled cross-checkpoint restart between AMReX and Paramesh. R\&D 100 Award, 2022.
  \item \textbf{MeshKit and Urban ECP.} PI for MeshKit (DOE NEAMS, 2009--2016): C++ toolkit for reactor mesh generation; parallel CoreGen produced a 101M-element MONJU mesh on 712 processors in about 7 minutes. Led Urban ECP: coupled WRF/HRRR $\rightarrow$ EnergyPlus $\rightarrow$ Nek5000 workflow with full 3D downtown Chicago domain. Published in \textit{Journal of Building Performance Simulation}, 2020.
\end{itemize}

\entryheader
  {University of Chicago}
  {2023 -- Present}
  {Staff At-Large}
  {Consortium for Advanced Science and Engineering (CASE), Chicago, IL}
\begin{itemize}
  \item Joint appointment supporting cancer pharmacogenomics and Earth science software collaborations; mentoring graduate students on software design and HPC workflows.
\end{itemize}

\entryheader
  {Arizona State University}
  {2007 -- 2009}
  {Research and Teaching Assistant}
  {Structural and Computational Mechanics Lab, Tempe, AZ}
\begin{itemize}
  \item FEM research on blast mitigation and design optimization; supported structural engineering coursework.
\end{itemize}

\entryheader
  {Earlier Experience}
  {2005 -- 2007}
  {Project Engineer, Wipro Technologies \textbar{} Engineering Intern, Tata Motors}
  {India}
\vspace{0.3em}

\section{Technical Skills}
\skillrow{Languages}{Python, C++, Fortran, Bash, R, SQL}
\skillrow{ML / AI}{PyTorch, TensorFlow, NumPy, Pandas, Xarray, Scikit-learn; Bayesian HPO (mlrMBO, Hyperopt/TPE), genetic algorithms (DEAP)}
\skillrow{AI Systems}{MCP server design (typed tools, scoped access, structured provenance), Claude API, LLM tool-use and agent workflows, local inference (llama.cpp, GGUF, quantization), Globus Compute for HPC-backed agents}
\skillrow{Correctness / Perf}{Numerical analysis (compensated summation, error-free transformations), roofline and bandwidth analysis, profiling, review of ported and machine-generated numerical code, reproducible benchmark design}
\skillrow{HPC / Systems}{MPI, OpenMP, HDF5, NetCDF, MOAB, Argobots, SZ3/ZFP, Intel XPU / CUDA, PBS/SLURM}
\skillrow{Workflow / DevOps}{GitHub Actions, Docker, Singularity, Parsl, Swift/T, release automation, CI/CD, Globus Compute}
\skillrow{Domains}{Climate and weather modeling, cancer pharmacogenomics, computational physics, reactor mesh generation}

\section{Selected Publications}
\begin{itemize}
  \item \textbf{Jain, R.} and Jacob, R. ``\href{https://agent4sc.github.io/}{Beyond Tool Execution: Evaluating Scientific MCP Interfaces with UXarray}.'' \textit{1st Workshop on Agentic AI for Large-scale Science (AGENT4SC) at IEEE eScience 2026}, Naples --- accepted, to be presented. Artifacts: \href{https://doi.org/10.5281/zenodo.21463232}{doi:10.5281/zenodo.21463232}.
  \item Partin, A., ..., \textbf{Jain, R.}, et al. ``\href{https://academic.oup.com/bib/article/27/1/bbaf667/7002013}{Benchmarking community drug response prediction models}.'' \textit{Briefings in Bioinformatics}, 2026.
  \item \textbf{Jain, R.}, Tang, H., Dhruv, A., and Byna, S. ``\href{https://doi.org/10.1109/SCW63240.2024.00043}{Enabling Data Reduction for Flash-X Simulations}.'' \textit{SC24 / DRBSD-10}, 2024.
  \item \textbf{Jain, R.}, Wozniak, J. M., et al. ``\href{https://web.cels.anl.gov/~woz/papers/IMPROVE_HPO_2024.pdf}{Cross-HPO: Optimizing Neural Networks for Cancer Drug Response}.'' \textit{CAFCW at SC24}, 2024.
  \item Wozniak, J. M., \textbf{Jain, R.}, et al. ``\href{https://bmcbioinformatics.biomedcentral.com/articles/10.1186/s12859-018-2508-4}{CANDLE/Supervisor: A workflow framework for machine learning applied to cancer research}.'' \textit{BMC Bioinformatics}, 2018.
\end{itemize}
Full list: \href{https://scholar.google.com/citations?user=bC77n9MAAAAJ}{Google Scholar}.

\section{Awards, Recognition, and Service}
\begin{itemize}
  \item Program committee member, 1st Workshop on Agentic AI for Large-scale Science (AGENT4SC) at IEEE eScience 2026.
  \item R\&D 100 Award, CANDLE (2023) \textbar{} R\&D 100 Award, FLASH-X (2022).
  \item DOE SBIR Phase I \& II (RGG reactor mesh generator; Kitware, 2014--2017).
  \item ATPESC Scholar, Argonne Training Program on Extreme-Scale Computing (2015).
  \item Best Paper Award, International Meshing Roundtable (2010).
\end{itemize}

\section{Education}
\textbf{University of Chicago} \hfill 2020 \\
M.S. in Computer Science

\vspace{4pt}
\textbf{Arizona State University} \hfill 2009 \\
M.S. in Structural Engineering \hfill Graduate Fellowship Recipient

\vspace{4pt}
\textbf{Indian Institute of Technology (ISM) Dhanbad} \hfill 2006 \\
B.Tech. in Mechanical Engineering

\end{document}
